\documentclass[gov]{harryopo-paper}

\title{关于开展办公文档智能化建设工作的通知}
\author{示例大学办公室\\ \fzfs\small 示例大学}
\date{2026年09月05日}


\begin{document}

\maketitle

\subsection*{一、总体要求}
\addcontentsline{toc}{subsection}{一、总体要求}

以提升公文生成效率与格式规范性为目标，全面推行模板驱动的文档生产方式。各部门要充分认识标准格式对公文权威性的重要意义，严格执行国家有关公文格式标准。

各学院、各部门应当结合自身实际，制定具体实施方案，确保各项工作落到实处、取得实效。

\subsection*{二、主要任务}
\addcontentsline{toc}{subsection}{二、主要任务}

\subsubsection*{（一）统一文档模板}
\addcontentsline{toc}{subsubsection}{（一）统一文档模板}

全校统一使用harryopo模板体系生成正式公文，涵盖通知、请示、报告、纪要等常用文种。

\subsubsection*{（二）规范生成流程}
\addcontentsline{toc}{subsubsection}{（二）规范生成流程}

\begin{enumerate}
  \item AI产出结构化Markdown中间态；
  \item 经办人核对内容后确认；
  \item 渲染引擎按GB/T 9704参数排版输出。
\end{enumerate}

\subsubsection*{（三）加强质量检查}
\addcontentsline{toc}{subsubsection}{（三）加强质量检查}

生成后应使用公文格式检查工具核对页边距、字体、字号与行距，确保与国家标准一致。

\subsection*{三、工作要求}
\addcontentsline{toc}{subsection}{三、工作要求}

（一）加强组织领导。各部门要明确责任人员，统筹推进本部门公文智能化工作。

（二）强化培训指导。信息中心要组织专题培训，帮助师生掌握工具使用方法。

（三）严格监督检查。学校办公室将定期抽查公文格式质量，通报检查结果。

\begin{quote}
注：本示例用于验证GB/T 9704公文模式（gov选项），正文三号仿宋、28磅行距、国标页边距。
\end{quote}

\end{document}